% Generated draft: requires mathematical and source review.
\documentclass[11pt]{article}
\usepackage[margin=25mm]{geometry}
\usepackage{fontspec}
\setmainfont{Latin Modern Roman}
\usepackage{amsmath,amssymb,mathtools,unicode-math}
\setmathfont{Latin Modern Math}
\usepackage{enumitem,needspace,xurl,hyperref}
\setlength{\parindent}{0pt}
\setlength{\parskip}{5pt}
\setlength{\emergencystretch}{3em}
\newcommand{\N}{\mathbb{N}}
\newcommand{\Z}{\mathbb{Z}}
\newcommand{\Q}{\mathbb{Q}}
\newcommand{\R}{\mathbb{R}}
\newcommand{\C}{\mathbb{C}}
\newcommand{\F}{\mathbb{F}}
\newcommand{\A}{\mathbb{A}}
\newcommand{\PP}{\mathbb{P}}
\newcommand{\E}{\mathbb{E}}
\newcommand{\eps}{\varepsilon}
\newcommand{\dd}{\,\mathrm{d}}
\begin{document}
\section*{UnsolvedMath ID: 30004389}
\textbf{Source problem code:} catalogue entry
\subsection{Definitions and mathematical statement}
Let \(M\) be a closed (compact without boundary), connected, smooth spin manifold of dimension \(d\geq 5\), and let \(G=\pi_1(M)\) be a finite abelian group of odd order. For a prime \(p\), call \(M\) \(p\)-atoral if, for every integer \(\ell\geq 1\) and every \(d\)-tuple \(c_1,\ldots,c_d\in H^1(M;\Z/p^\ell)\),
\[
\bigl\langle c_1\smile\cdots\smile c_d,[M]_{p^\ell}\bigr\rangle=0\quad\text{in }\Z/p^\ell,
\]
where \([M]_{p^\ell}\in H_d(M;\Z/p^\ell)\) is the reduction of the fundamental class determined by the spin orientation. Assume this condition for every prime \(p\mid |G|\). If \(G\) is trivial, this collection of conditions is vacuous.

Using the standard real convention, write
\[
\alpha(M)=\operatorname{Ind}_{C^*_{\R}G}\bigl(D_M\otimes\mathcal L_G\bigr)\in KO_d(C^*_{\R}G)
\]
for the Rosenberg invariant: the real higher index of the spin Dirac operator \(D_M\), twisted by the flat bundle \(\mathcal L_G=\widetilde M\times_G C^*_{\R}G\). Here \(C^*_{\R}G\) is the real group \(C^*\)-algebra; its maximal and reduced completions agree for finite \(G\).

The proposed statement is that, for every such spin manifold,
\[
\alpha(M)=0\quad\Longrightarrow\quad
\text{there exists a smooth Riemannian metric }g\text{ on }M
\text{ with }\operatorname{Scal}_g(x)>0\text{ for every }x\in M.
\]
The supplied original formulation asks whether a non-spin positive-scalar-curvature existence theorem extends to this spin setting. The supplied cleaned statement identifies the invariant as the Rosenberg invariant; the real convention above makes that interpretation explicit, although the record does not specify its coefficient convention. This is a proposed implication requiring mathematical review, not an asserted theorem.
\subsection{Short English statement}
Does every closed connected smooth spin manifold of dimension at least five, with finite abelian fundamental group of odd order and the stated vanishing of all top-degree products of degree-one classes over \(\Z/p^\ell\) for each relevant prime and every \(\ell\geq1\), admit a metric of everywhere positive scalar curvature whenever its Rosenberg invariant vanishes? The supplied record labels this spin extension open.
\subsection{Sources}
\begin{itemize}
\item Local UnsolvedMath registry, permanent ID \(30004389\), source problem code \texttt{catalogue entry}, titled ``Positive Scalar Curvature for Spin Manifolds with Odd Fundamental Group,'' in the Oberwolfach Reports Open Problems collection. No individual proposer is supplied; the proposed year is recorded as 2020.
\item Supplied primary citation: \emph{Manifolds and Groups (2021)}, Oberwolfach Reports, pp.~417--463, DOI \texttt{10.4171/owr/2020/7}. \url{https://doi.org/10.4171/owr/2020/7}. The record says the official report was checked and labels the statement \texttt{corrected\_verified}; these are supplied verification claims, not independent verification for this draft.
\item Supplied literature source labeled ``Hanke 2021'': \url{https://doi.org/10.2140/gt.2021.25.497}. No article title or precise theorem locator is supplied.
\item Literature assessment dated exactly \texttt{2026-08-22}: the supplied curation labels the problem open, describes the spin extension with vanishing alpha-invariant as conjectural, and identifies connective K-homology as the missing method. This assessment is supplied dated curation derived from the Oberwolfach material, not independently verified here; its time-sensitive status and cited claims require independent checking.
\end{itemize}
\end{document}
